\section{Conclusão}
O estudo demonstrou a viabilidade de implementar um modelo de previsão fotovoltaica que combina informações meteorológicas, histórico de geração e parâmetros físicos dos equipamentos. O LSTM direto apresentou o menor erro no teste simulado, enquanto a RNN liderou a validação e o Random Forest foi competitivo, especialmente sob chuva e nebulosidade. A aprendizagem residual melhorou a estimativa física, mas não foi a abordagem mais precisa do comparativo.

Conclui-se que a seleção das variáveis, a coerência física e a avaliação dos erros são essenciais para o desenvolvimento de um projeto de Machine Learning aplicado à energia solar. A arquitetura simples é adequada como prova de conceito acadêmica. A precisão operacional e a capacidade de generalização ainda precisam ser demonstradas com medições reais.